\chapter{Theoretical foundation and bibliographical documentation}
\label{cap:cap2}

\section{Relevant Concepts and Theories}

\subsection{Physiological Caloric and Macro-nutrient Modeling}
Before a software architecture can automate meal generation, it must establish a baseline 
understanding of human energy metabolism. The foundation of personalized nutrition relies 
on two primary metrics: Basal Metabolic Rate (BMR) and Total Daily Energy Expenditure 
(TDEE).

BMR represents the minimal amount of energy, measured in calories, that the human body requires 
to maintain vital vegetative functions—such as cell production, respiration, and temperature 
regulation—at absolute rest in a post-absorptive state \cite[Ch.~2]{techreport:fao:energy}. 
Because BMR accounts for roughly 60\% to 75\% of an individual's daily energy output, it serves 
as the mandatory baseline for any personalized coaching algorithm.

TDEE, on the other hand, represents the total number of calories an individual burns 
over a 24-hour window. It builds upon the BMR by incorporating additional metabolic 
demands, primarily the energy expended during physical exercise, daily functional movement, 
and the thermic effect of food processing, all of which dynamically scale an individual's 
baseline energy requirement based on adult activity vectors \cite[Ch.~5]{techreport:fao:energy}.

To establish a precise nutritional target, the platform must dynamically calculate an 
individual's Total Daily Energy Expenditure (TDEE). This process is executed in a 
two-stage mathematical pipeline: first, determining the Basal Metabolic Rate (BMR), 
and second, applying an activity multiplier coefficient based on the user's lifestyle 
density.

The baseline metabolic requirements are computed using the Mifflin-St Jeor 
equation \cite{Mifflin1990}, which models resting energy expenditure based on 
weight ($w$ in kg), height ($h$ in cm), age ($a$ in years), and a discrete biological sex 
constant ($s$). The mathematical formulation is defined as:

\begin{equation}
BMR = 10w + 6.25h - 5a + s
\end{equation}

where the gender-specific parameter $s$ is defined conditionally as:

\begin{equation}
s = \begin{cases} 
5, & \text{if user is male} \\
-161, & \text{if user is female}
\end{cases}
\end{equation}

To scale the resting metabolic rate to a realistic daily caloric maintenance target, 
the calculated BMR is mapped against standard metabolic activity multipliers derived 
from historical evaluations of the Harris-Benedict principle \cite{Roza1984}. The 
platform defines the final TDEE using an activity factor ($AF$) corresponding to five 
distinct physiological movement groups:

\begin{equation}
TDEE = BMR \times AF
\end{equation}

The activity factor coefficient $AF$ is assigned according to the user's weekly exercise 
intensity vectors:
\begin{align*}
AF &= 1.200 \quad \text{(Inactive: little to no exercise)} \\
AF &= 1.375 \quad \text{(Low: light exercise 1--3 times/week)} \\
AF &= 1.550 \quad \text{(Medium: moderate exercise 3--5 times/week)} \\
AF &= 1.725 \quad \text{(High: intense exercise 6--7 times/week)} \\
AF &= 1.900 \quad \text{(Intense: daily heavy physical labor or athletic sports)}
\end{align*}

\subsection{Multi-Layer Perceptron (MLP) Classifier Foundations}

To automate the selection of recipe baselines within the system, a Multi-Layer Perceptron 
(MLP) architecture is utilized as an intelligent heuristic filter. Instead of executing 
resource-intensive database queries to scan thousands of potential meal profiles, the 
network processes the primary user macro-nutrient requirements instantly to isolate a 
highly ranked subset of candidate configurations.

\subsubsection{Structural Topology and Layer Configurations}
The neural network is structured as a feedforward network containing an input layer, two 
fully connected hidden layers, and an unnormalized dense output layer. The input features 
are mapped as a 3-dimensional vector containing the targets for calories, protein, and the 
categorical meal identifier:

\[
x = \begin{bmatrix} \text{Calories} \\ \text{Protein} \\ \text{Meal\_Type\_ID} \end{bmatrix}
\]

The data propagates through the network via successive linear matrix transformations. At each 
layer, the incoming vector is multiplied by a weight matrix ($W$) and combined with a bias 
vector ($b$):

\begin{equation}
\text{Output} = W \cdot x + b
\end{equation}

To introduce non-linear mapping capabilities, the hidden layers apply the Rectified Linear 
Unit (ReLU) activation function, which truncates negative values to zero. The network 
configuration expands the 3-dimensional input into a primary hidden layer of 512 neurons, 
compresses the features into a secondary hidden layer of 256 neurons, and finally projects 
onto the output layer.

\subsubsection{Top-K Output Selection}
The final layer contains a total number of nodes equal to the total number of pre-compiled 
recipe variations in the system registry. Rather than computing an expensive probability 
distribution across all possibilities, the network yields raw structural evaluation scores 
known as logits. 

The application logic extracts the highest-scoring indices from this output vector using a 
Top-K maximum selection operator. This isolates a focused array of candidate recipe IDs 
(where $k=5$ for standard targets or $k=15$ for simple complexity structures) to be passed 
downstream for exact ingredient weight optimization.

\subsection{Linear Programming and Constraints Optimization}

To resolve the exact ingredient quantities necessary to fulfill a user's 
macro-nutrient targets, the system utilizes Linear Programming (LP). While 
the Multi-Layer Perceptron architecture acts as a heuristic filter to select 
the most structurally appropriate recipe templates, the Linear Programming solver 
operates as an exact mathematical algorithm. It determines the optimal allocation 
of continuous variables (ingredient weights in grams) under strict linear boundary limits.

\subsubsection{Problem Formulation and the Objective Function}
The optimization process treats each ingredient within the selected recipe as a decision 
variable $x_i$, representing the mass of that ingredient in grams. The primary objective 
is to minimize a linear cost or deviation function, ensuring the meal remains as close as 
possible to an ideal baseline structure while satisfying the user's daily goals. The 
general mathematical formulation minimizes the objective function:

\begin{equation}
\min \sum_{i=1}^{n} c_i x_i
\end{equation}

Where $n$ is the total number of ingredients in the candidate meal, and $c_i$ represents 
the constant coefficients dictating the optimization priority or scaling factors for 
each specific ingredient.

\subsubsection{Linear Constraints and Boundary Conditions}
The solution space is strictly bounded by a system of linear inequalities derived from 
the user's nutritional parameters. For a target plan requiring specific calorie 
bounds ($E$), a minimum protein threshold ($P$), and a minimum fat threshold ($F$), 
the system maps out the continuous constraints as follows:

\begin{equation}
\sum_{i=1}^{n} e_i x_i \le E_{\text{target}}
\end{equation}

\begin{equation}
\sum_{i=1}^{n} p_i x_i \ge P_{\text{target}}
\end{equation}

\begin{equation}
\sum_{i=1}^{n} f_i x_i \ge F_{\text{target}}
\end{equation}

Where $e_i$, $p_i$, and $f_i$ represent the constant nutritional values of 
ingredient $i$ per single gram (for calories, protein, and fat respectively). 
Additionally, non-negativity and structural boundaries are applied to prevent 
negative food weights and ensure realistic ingredient portions:

\begin{equation}
x_{\min, i} \le x_i \le x_{\max, i}, \quad \forall i \in \{1, \dots, n\}
\end{equation}

By evaluating these bounded intersections, the solver instantly locates the precise 
feasible region, bypassing manual adjustment loops to output a mathematically sound 
meal composition.

\section{Comparative Analysis of Existing Solutions}

The digital nutrition ecosystem features various architectural approaches to 
software-assisted macro-nutrient management. These solutions generally fall into 
two categories: manual logging utilities and rigid algorithmic recipe generators. 

\begin{itemize}
    \item \textbf{Manual Logging Utilities (e.g., MyFitnessPal):} These platforms operate 
    primarily as relational database management interfaces. Users manually search, select, 
    and log food entries post-consumption \cite{url:myfitnesspal}. Caloric guidelines are treated as static targets 
    derived from baseline formulas, placing the operational burden of macro-nutrient 
    mathematical alignment entirely on the end user.
    \item \textbf{Dynamic Metabolic Trackers (e.g., MacroFactor):} These systems use advanced 
    expenditure analytics to continuously recalculate metabolic expenditure based on daily 
    weight and logging inputs \cite{url:macrofactor}. While highly accurate from a physiological modeling 
    perspective, they remain bound to manual data-entry paradigms and lack automated meal 
    generation capabilities.
    \item \textbf{Automated Rule-Based Planners (e.g., Eat This Much):} 
    Standard automated meal planners utilize deterministic decision trees or 
    static search queries to select pre-authored recipes from a fixed 
    index \cite{url:eatthismuch}. These platforms lack micro-granular ingredient 
    flexibility; they cannot dynamically scale individual ingredient quantities 
    to satisfy strict mathematical constraint configurations across an entire day.
\end{itemize}

A technical comparison of these architectural paradigms relative to the proposed system 
is outlined in Table~\ref{tab:comparative_analysis}.

\begin{table}[!htbp]
\centering
\caption{Comparative Matrix of Existing Digital Nutrition Solutions}
\label{tab:comparative_analysis}
\begin{tabular}{|l|c|c|c|c|}
\hline
\textbf{Technical Criterion} & \textbf{MyFitnessPal} & \textbf{MacroFactor} & \textbf{Standard Planners} & \textbf{Proposed System} \\ \hline
Automated Generation        & No                    & No                   & Yes                        & \textbf{Yes}             \\ \hline
Heuristic Prep-Filtering    & No                    & No                   & No                         & \textbf{Yes (MLP)}       \\ \hline
Exact Gram-Level Scaling    & No                    & No                   & No                         & \textbf{Yes (LP)}        \\ \hline
Dynamic Macro Adjustments   & No                    & Yes                  & No                         & \textbf{Yes}             \\ \hline
Minimal Operational Burden  & No                    & No                   & Moderate                   & \textbf{High}            \\ \hline
\end{tabular}
\end{table}

\section{Identification of Gaps in Current Solutions}

The comparative analysis highlights three primary architectural gaps within the current 
state of the art:

\begin{enumerate}
    \item \textbf{High User Cognitive Load:} Traditional systems depend on continuous 
    manual interaction loops, requiring the user to manually weigh, search, and log 
    individual ingredients post-consumption. This tedious operational cycle quickly 
    induces tracking fatigue, which serves as the primary driver for high user abandonment 
    rates in digital dietetics software.
    \item \textbf{Rigidity of Menu Databases:} Existing automated planners treat recipe 
    structures as static lookups. If an isolated recipe configuration fails to meet a 
    user's calculated macro-nutrient targets exactly, the system cannot modify its internal 
    components. It must either reject the recipe entirely or present an unaligned, 
    mathematically inaccurate menu choice to the user.
    \item \textbf{Lack of Multi-Stage Architecture:} Current software systems treat 
    automated meal planning either as a pure database search or a pure mathematical 
    constraint optimization problem. Applying complex optimization models globally 
    across thousands of raw database rows fails to scale efficiently on standard backend 
    web servers, resulting in substantial processing latencies that degrade the user 
    experience.
\end{enumerate}

\section{Expected System Specifications}

To resolve the identified limitations in existing digital 
dietetics platforms, the proposed system is engineered around a 
derived set of technical and functional specifications:

\begin{itemize}
    \item \textbf{Functional Autonomy:} The platform must automatically synthesize 
    a complete, cohesive daily meal blueprint using calculated target calories, 
    protein thresholds, and minimum fat boundaries. This removes the necessity of 
    manual food combination planning and reduces the user's operational interaction 
    loop to a simple configuration step.
    \item \textbf{Hybrid Pipeline Execution:} The system backend must execute a two-stage 
    processing pipeline to combine performance speed with mathematical precision. 
    A machine learning classification layer (MLP) must handle initial structural baseline 
    filtering from the database, which a dedicated constraint solver (Linear Programming) 
    then scales down to individual ingredient masses in grams.
    \item \textbf{User-Defined Complexity Constraints:} The menu generation pipeline must 
    dynamically adapt candidate food pools based on specified user lifestyle parameters. 
    The system must safely isolate simple, time-efficient recipes containing minimal 
    preparation steps from standard multi-ingredient variants, without causing mathematical 
    failures or violating nutritional safety margins.
\end{itemize}